\section{Task Suite Data Details}
\label{app:data-details}


\providecommand{\ctc}[1]{\ensuremath{\mathcal{O}_T(#1)}}
% A query, then the answer it should produce.
\providecommand{\ex}[2]{#1~$\rightarrow$~\texttt{#2}}
% A class block header spanning the table.
\providecommand{\exblk}[1]{\multicolumn{2}{@{}l}{#1} \\}

\begin{table}[t]
\centering
\scriptsize
\setlength{\tabcolsep}{4pt}
\renewcommand{\arraystretch}{1.12}

\begin{tabular}{@{}
  >{\raggedright\arraybackslash}p{0.20\linewidth}
  >{\raggedright\arraybackslash}p{0.78\linewidth}
@{}}
\toprule
\textbf{Dataset} & \textbf{Example} \\
\midrule
\exblk{\textbf{\(\mathbf{O}(N)\): One-pass tasks}}
\midrule

NQ
  & \ex{\emph{``who sold out jesus for 30 pieces of silver''}}{[1]} \\
HotpotQA (bridge)
  & \ex{\emph{``What company did Rex Maughan aquire?''}}{[8, 9]} \\
NIAH-contra
  & \ex{\emph{``$\dots$a hypertonic solution of 14.4\% has an \textbf{ir}reversible ciliostatic
    effect''}}{[20]}, the one claim saying \emph{reversible} \\
BEIR SciFact
  & \ex{\emph{``0-dimensional biomaterials show inductive properties.''}}{[1]} \\
BEIR FiQA
  & \ex{\emph{``Where should I park my rainy-day / emergency fund?''}}{[2, 8, 10, 12, 19]} \\
MS MARCO
  & \ex{\emph{``what is the difference between fables and folktales? for kids''}}{[2]} \\
MS MARCO rerank
  & the same 20-passage pool, every passage scored $\rightarrow$
    \texttt{Ranking: [2], [1], [3], \dots} \\
OBLIQ
  & \ex{\emph{``$\dots$written in a similar style and voice to: The slowness and calmness of the
    lodge came with genuine inquisitiveness$\dots$''}}{[14, 23, 29]} \\
Oolong
  & 10 labelled sentence pairs; \ex{\emph{``which of the labels is the most
    common?''}}{Label: incorrect} \\
Outlier (Amazon)
  & 20 reviews, 16 rated 4-star and 4 rated 5-star $\rightarrow$ \texttt{Outliers: 4, 7, 11, 17} \\
Outlier (wiki, fixed \(M\))
  & 57 chunks from 3 articles (43 / 11 / \textbf{3}) $\rightarrow$ \texttt{Outliers: 5, 17, 54} \\
Absence
  & 90 numbered sentences, 3 deleted from the copy $\rightarrow$ \texttt{["The interior is what",
    "In opposite chapels are", \dots]} \\

\midrule
\exblk{\textbf{Higher-complexity tasks}}
\midrule

Outlier (wiki, scale-\(k\))
  & 55 chunks from 10 articles (9 / 7 / 7 / $\dots$ / \textbf{3}) $\rightarrow$
    \texttt{Outliers: 26, 41, 42} \\
Grouping (OpenAlex)
  & \ex{\emph{``Cluster these 20 papers into 15 groups.''}}{\{"groups": [\{"doc\_ids":
    [4, 5, 12]\}, \dots]\}} \\
Contradiction
  & \texttt{[2]} \emph{``$\dots$has a reversible ciliostatic effect''} vs.\ \texttt{[12]}
    \emph{``$\dots$has an irreversible ciliostatic effect''} $\rightarrow$
    \texttt{[[2,12], [8,11], [13,17]]} \\
xabsence
  & 19 claims, 8 matched pairs, 3 with no twin in the other corpus $\rightarrow$
    \texttt{Unmatched: 1, 3, 14} \\
strmatch
  & \texttt{[1]} and \texttt{[11]} share \emph{``kossuth westward taco''} $\rightarrow$
    \texttt{[[1,11], [2,8], [15,19]]} \\
qdmatch (NQ)
  & \texttt{[7]} \emph{``when did 10 shilling note go out of circulation''} $\leftrightarrow$
    \texttt{[29]} $\rightarrow$ \texttt{[[7,29], [14,28], [19,40]]} \\
qdmatch (HPQA)
  & \texttt{[8]} \emph{``\,`Lost!' is a song by a British rock band formed in what year?''}
    $\leftrightarrow$ \texttt{[25]} and \texttt{[34]} \\
qdmatch (FiQA)
  & \texttt{[25]} \emph{``Pay off credit card debt or earn employer 401(k) match?''}
    $\leftrightarrow$ \texttt{[55]} \\
Reorder
  & 50 shuffled segments of \emph{Flatland} $\rightarrow$ \texttt{[31, 19, 12, 22, 37, \dots]} \\
Textgroups
  & noun counts $34{+}31{+}5 = 70$ and $13{+}17{+}40 = 70$ $\rightarrow$
    \texttt{[[1,11,12], [2,14,17]]} \\

\bottomrule
\end{tabular}

\caption{\textbf{One real example per task.} All content is verbatim from the evaluation sets,
elided with \texttt{\dots} where a passage is too long to show; document IDs are 1-indexed exactly
as the model sees them. Appendix~\ref{app:examples} shows each example in full, with the task
instruction, the surrounding corpus and the distractor types.
}
\label{tab:task-examples}
\end{table}

\prasann{TODO this is a starter Claude generation from suite docs, need to add more detail / clean up}
% The CoT ablation in Table~\ref{tab:cot-ablation} uses a smaller-scale variant of the outlier task built from the Amazon Reviews 2023 corpus \citep{Hou2024BridgingLA}, distinct from the wiki-100w outlier data used for the cross-task scaling figure (Figure~\ref{fig:sec5-cross-task}).

% \paragraph{Construction.} Each example contains $N{=}30$ short review documents (review title + body, no rating or category metadata visible to the model). Examples are one of two types, mixed 50/50 in the dataset via the \texttt{--rating-ratio} flag of \texttt{generate\_review\_outlier\_data.py}:

% \begin{itemize}
% \item \textbf{Rating outlier:} a majority of reviews share one star rating (1--5) and 3 outlier reviews carry a different rating. Query: \emph{``Can you find outliers with the least common rating in this data?''}
% \item \textbf{Category outlier:} a majority of reviews come from one product category and 3 outliers come from a different category. Categories are drawn from the HF config names of the Amazon-Reviews-2023 dump (e.g.\ \texttt{Books}, \texttt{Electronics}, \texttt{Beauty\_and\_Personal\_Care}, \texttt{Home\_and\_Kitchen}). Query: \emph{``Can you find outliers from the least common category?''}
% \end{itemize}

% Ground truth is derived from the structured metadata (star rating field, HF config name) and \emph{not} from any LLM labeling---the model must read the review content to identify outliers because the visible text strips both fields.

% \paragraph{Splits used in Table~\ref{tab:cot-ablation}.} 5{,}000 train examples, 500 eval examples, all with $N{=}30$, 3 outliers each. CoT targets are computed deterministically from the same metadata: the \emph{faithful} target names the majority and minority labels (e.g.\ ``Most reviews are 4-star ratings and the outliers are 5-star reviews.''); \emph{repeat-IDs} emits the gold outlier doc IDs before the answer; \emph{mislabeled} uses a randomly-permuted (wrong) majority/minority pair while keeping the template; \emph{dummy} is a content-less template phrase; \emph{random words} samples 50 content words from the input documents. All targets share the same answer format (a comma-separated list of outlier doc IDs).

% \paragraph{Why a separate dataset.} The wiki-100w outlier task involves grouping 100-word Wikipedia chunks by their source article, which gives the model a strong topical-similarity prior. The Amazon Reviews variant intentionally weakens that prior---ratings and categories often produce reviews of similar surface style---which makes the task harder to learn without natural-language guidance, and making the CoT ablation more informative.


\begin{claudeblock}
% Full-width [t] floats: 10 and 12 panels do not fit a column, and both are reference figures the
% reader jumps to rather than reads in sequence.
\begin{figure}[t]
\centering
\includegraphics[width=\textwidth]{figures/ctc_suite_high.pdf}
\caption{\textbf{Every high-CTC task, full versus chunked attention.} The ten $O_T(NM)$,
$O_T(N^2)$ and $O_T(N^3)$ tasks of the suite, Qwen3.5-4B, solid = full attention and dashed =
document-chunked with mask mixing. The panels share one axis so the \emph{gaps} are comparable by
eye; the absolute scores are not, since each task uses its own metric. Re-ordering has no 32k rung
(never evaluated, not a measured zero). These are the individual tasks behind the high-CTC average
in Figure~\ref{fig:multitask_fig4}.}
\label{fig:suite-high}
\end{figure}

\begin{figure}[t]
\centering
\includegraphics[width=\textwidth]{figures/ctc_suite_low.pdf}
\caption{\textbf{Every low-CTC task, full versus chunked attention.} The twelve $O_T(N)$ tasks,
drawn on the same axes and with the same conventions as Figure~\ref{fig:suite-high}. The two arms
are close to indistinguishable on almost every panel --- which is the contrast the paper rests on:
the same mask, the same models and the same ladders cost little here and a great deal in
Figure~\ref{fig:suite-high}. Outlier fix-$k$ and absence have no 32k rung.}
\label{fig:suite-low}
\end{figure}
\end{claudeblock}
\begin{claudeblock}

\subsection{Shared construction protocol}
\label{app:construction}

Two things are common to every task in the suite, and both matter for reproducing the numbers.

\paragraph{The length ladders vary only the distractors.} Each task is evaluated at five context
budgets (2k, 4k, 8k, 16k, 32k tokens). These are \emph{not} five independently sampled datasets.
For each task we fix one canonical question set --- the largest existing rung, which already carries
every gold document and the deepest distractor pool --- and derive every shorter rung from it by a
nested prefix: the gold documents, the query and the answer are held byte-identical, a single
deterministic shuffled order is imposed on the distractor pool, and rung $n$ keeps the golds plus
the first $n - |\text{gold}|$ distractors. The combined list is then shuffled and every index field
(gold document indices, hard-negative indices, and any answer string derived from position) is
remapped.

The consequence is that a shorter rung's distractor set is a strict subset of a longer rung's, so
context length is the only variable that changes along a ladder and rung-to-rung comparisons are
not confounded by question difficulty. Every rung holds at least 500 examples. Where a task's
documents are too long for a budget to be met by document count alone --- OBLIQ, whose median
example has seven long multi-gold posts --- the rung is built by a gold-preserving per-example
token budget instead, and its 2k rung is omitted rather than truncated, because truncating it
silently drops gold documents and scores the task at zero.

\paragraph{Training data is i.i.d.\ with evaluation but disjoint from it.} Unless stated otherwise,
each task is trained for one epoch on 20{,}000 examples split evenly across the same five context
budgets used at evaluation, so the training distribution matches the ladder rather than a single
length. Training examples are drawn from the same generator as the evaluation examples for that
task, with source documents held disjoint --- fillers are sampled from abstracts, articles or
passages that supply no gold in that example, and, for the tasks with a published split, the
benchmark's own train/test division is respected. OBLIQ is the sole exception: it has no usable
training split and its training data is synthesized, under a clean-room guard described above.
\end{claudeblock}

\subsection{Tasks: Overview}

\subsubsection{\texorpdfstring{\ctc{N}}{OT(N)} tasks: one pass over the corpus}\autonote{UNCITED ON PURPOSE: BEIR, SciFact, FiQA, MS MARCO and the ms-marco cross-encoder are named in prose with no \texttt{\textbackslash citep} --- these are the spots that want one, since only you add citations. GovReport was in this list and has been dropped with the NarrativeQA/GovReport paragraph. The NQ and HotpotQA paragraphs are edited versions of the old blurbs --- NQ now records that current data caps hard negatives at $\sim$10\% via the CE filter, which the old text did not.}

These tasks decompose into per-document oracle calls $\mathrm{isrelevant}(d_i, q)$: a single sweep suffices, and nothing about the answer depends on how two fillers relate to each other.

\textbf{Factoid QA (NQ).} \emph{Gold-ID F1.} NaturalQuestions-open \citep{izacardNQ}, with all documents --- gold, hard negative and filler --- drawn from the 100-word Wikipedia DPR corpus \citep{Karpukhin2020DensePR} served by pyserini \citep{Lin2021PyseriniAP}, so every document shares one surface format. For each question we BM25-search \citep{Robertson1976RelevanceWOBM25} for a passage containing an answer string (that becomes the gold), take top BM25 hits that do \emph{not} contain the answer as hard negatives, and fill the remaining slots with random corpus passages. Recent versions of this data hold hard negatives to $\sim$10\% of the pool with the cross-encoder filter described above; earlier all-hard-negative pools made the task an adversarial ranking problem rather than a retrieval one.

\textbf{Multi-hop QA (HotpotQA).} \emph{Gold-ID F1.} The bridge subset of HotpotQA \citep{Yang2018HotpotQAAD}, where answering requires composing two gold passages; the corpus is again wikipedia-100w with BM25 hard negatives. Because the answer needs a fixed (two) number of retrieval steps, this remains \ctc{N}.
\claude{The suite runs use no chain of thought for this task: the target is the gold document ID
set and the metric is gold-ID F1, so the model retrieves rather than answers.}
% ⚠ CORRECTED: this said "We include a short natural-language chain of thought generated with
% gemini-2.5-flash, grounded in the gold paragraphs and the answer." The current HotpotQA ladder
% carries no CoT field at all (eval JSONL keys are answers / documents / gold_doc_indices /
% hard_neg_indices / queries / source) and every harvested grade JSON has cot_mode unset. The
% gemini CoT belongs to the earlier per-N HotpotQA data, not the suite ladders.

\textbf{Needle contradiction (NIAH-contra).} \emph{Set-F1.} An \ctc{N} control derived from the \ctc{N^2} contradiction task, holding the substrate fixed while dropping the complexity class: one member of a gold contradiction pair is promoted to the query (``\emph{find the document that contradicts this claim}'') and the other stays in the corpus as the needle. Every other claim, including the other gold pairs, is a distractor. Comparing this against contradiction isolates the effect of all-pairs search from the effect of PubMed claim text.

\textbf{BEIR SciFact / FiQA.} \emph{Set-F1.} Two BEIR tasks framed in-context: SciFact (scientific claim $\to$ abstract, single-gold) and FiQA (financial opinion QA, median 2 golds/query, 648 queries). Both mine distractors by BM25 from the dataset's own corpus. FiQA is sparsely judged, so plain BM25 distractors are riddled with unlabeled positives; we score (query, candidate) pairs with a cross-encoder and keep a candidate only if it scores at least a margin below the gold, then fill the rest of the pool with random real corpus documents.

\textbf{MS MARCO retrieval and reranking.} \emph{Set-F1; MRR@10 (NDCG@10 for the HELMET-format variant).} Passages come from the \texttt{msmarco-v1-passage} index. Hard negatives are taken from SentenceTransformers' precomputed BM25+dense mining and filtered by a margin on precomputed cross-encoder scores, so removing false negatives requires no model run; random fill is sampled as random pids from the 8.8M collection, which are byte-identical in format to gold and hard passages. The rerank variant reuses the identical pool but asks for the documents in ranked order, with every document carrying a real cross-encoder relevance score --- so the target order is defined for the whole pool, not just the judged part.

\textbf{OBLIQ retrieval.} \emph{Set-F1.} OBLIQ-Bench (\texttt{dianetc/OBLIQ-Bench}) poses subjective, long-form queries (``\emph{find a document written in a similar style to this passage}'') over its own corpora; we port it in-context by forcing all qrel positives into the example and filling with BM25-mined distractors from the same subset. Its queries are long and its documents are full posts, making it the longest-per-document retrieval task in the suite.

\begin{claudeblock}
\emph{Synthetic training data.} OBLIQ is the one task with no usable training split: the benchmark
ships only a few hundred real queries per subset, which yields fewer than 2{,}000 in-context
training examples and leaves the model badly undertrained (in an early run, roughly 70\% of its
predictions at $N{=}30$ were one of two canned strings). The bottleneck is \emph{query} poverty, not
document poverty --- each subset corpus holds thousands of documents --- so we synthesize only the
scarce element and keep the real corpus. For each seed document $D$:

\begin{enumerate}
\item An LLM (Qwen3-8B) writes an oblique query $Q$ that $D$ satisfies, few-shot primed with real
(document, query) pairs so the synthetic queries match the benchmark's register --- the real queries
turn out to be multi-gold descriptions of a property or theme, not ``find the one document that
$\dots$''.
\item An \emph{obliqueness} filter rejects any $Q$ with word-Jaccard overlap against $D$ above a
per-subset cap (0.55 for the writing subset, 0.35 elsewhere), so a keyword shortcut cannot separate
the golds --- matching real OBLIQ, whose queries are deliberately non-lexical.
\item BM25 retrieves candidate documents for $Q$ from the same subset.
\item A \emph{multi-gold judge} asks, for each top candidate, whether it \emph{also} satisfies $Q$.
A yes becomes an additional gold; a no becomes a hard distractor. This is the step that matters: it
both matches the multi-gold eval distribution and removes false negatives by labelling them rather
than discarding them.
\item The judged example is emitted through the same unified-format path as the real data.
\end{enumerate}

Two properties are worth stating explicitly. \textbf{Only training data is synthetic} --- the
held-out evaluation set remains the real 488-example OBLIQ test set, since synthesizing the
evaluation would confound the capability question with generator bias. And the generator runs
\textbf{clean-room}: every document appearing anywhere in the evaluation file, as a gold
\emph{or} as a distractor, is removed from the corpus before any seeding or BM25 retrieval, so no
evaluation document can enter training as a seed, a gold or a distractor. Only the few-shot
\emph{query} exemplars come from the evaluation set.

Scaling this to roughly 15{,}000 examples improved a matched 0.8B run over the $\sim$1{,}800-example
real-query baseline on the real test set (gold-ID F1 $0.317 \to 0.339$, exact match $0.049 \to
0.090$). The F1 gain is about one standard error at this evaluation size and is suggestive rather
than conclusive; the exact-match gain is larger. Notably, more data does not lift OBLIQ off its
$\sim$0.33 floor, which we read as the task being genuinely hard rather than merely
data-starved.
\end{claudeblock}

\textbf{Oolong.} \emph{Partial-credit score.} A port of Oolong \citep{Bertsch2025OolongEL}: a long context of labeled classification items, each tagged with a date and a user id, and a distributional question --- most/least common label, most frequent user, when label A overtook label B. We use the \emph{unlabeled} context variant, so the model must classify each item and then aggregate. Because the label set is fixed and each item is judged independently before a single aggregation step, this is \ctc{N} --- which is exactly what makes it a useful contrast with our outlier tasks below, which look superficially similar but have an unbounded category set.

% ⚠ REMOVED FROM THE SUITE: NarrativeQA and GovReport are not among the 22 roster tasks
% (they are single-document HELMET ports with non-ID metrics). Commented, not deleted.
% \textbf{NarrativeQA and GovReport (HELMET ports).} \emph{Token-F1; ROUGE.} Two HELMET \citep{Yen2024HELMETHT} long-document tasks kept in their native form: question answering over a single book or screenplay \citep{Kocisk2017TheNR}, and summarization of a government report. Both are single-document tasks --- there is no cross-document structure to exploit --- and serve as the ``long context is not the same as many documents'' anchor of the suite. These are the two tasks whose metrics are not ID-based, and we do not use them for headline gap-growth claims.

\textbf{Absence.} \emph{Set-F1 over removed IDs.} The inverse of needle-in-a-haystack, following AbsenceBench (arXiv:2506.11440): the model sees the full numbered corpus, then a second copy with some elements deleted, and must name what is missing. We derive it over our own corpora: each document is deleted independently with probability $p$, scattered single deletions being the hard regime. Establishing that item $i$ is absent requires checking it against the second corpus, hence \ctc{N^2}.


\subsubsection{\texorpdfstring{\ctc{NM}}{OT(NM)} tasks: categorization}\autonote{Both remaining tasks are edited old blurbs. Two claims changed on purpose: the old text
said the wiki outlier task ``resembles OOLONG'', which read as if it were the same task --- it is now
stated as the contrast (open vs.\ fixed category set) --- and the scale-$k$ vs.\ fixed-$M$ control is
named explicitly. Grouping now describes the labeled variant, since the unlabeled one was cut.}

Here the oracle operation compares a document against a \emph{category} rather than a query, and the number of categories $M$ is not fixed by the task definition --- it is a property of the particular corpus, which the model must itself infer. A naive algorithm compares each document against each discovered category, giving \ctc{NM}.

\textbf{Outlier detection (Wikipedia).} \emph{Set-F1.} Each context is filled with 100-word chunks from a handful of Wikipedia articles with imbalanced counts; the answer is the ID list of chunks belonging to the article with the fewest chunks. Article titles are withheld from the model (they would trivially leak the outlier), so the grouping signal has to come from the body text. We run two variants: \emph{scale-$k$}, where the number of source articles grows with $N$ (so $M$ grows, and the task is genuinely \ctc{NM}), and a \emph{fixed-$M$} control where the article count stays constant as $N$ grows --- the control that separates ``the gap widens because of $N$'' from ``the gap widens because of $M$''. This differs from Oolong precisely in that the category set is open: there is no fixed label vocabulary to aggregate over.

\begin{claudeblock}
\textbf{Outlier detection (Amazon Reviews).} \emph{Set-F1.} The same outlier structure on a
different substrate: each example is $N{=}30$ customer reviews from the Amazon Reviews 2023 corpus,
shown as review title and body only. Examples are mixed 50/50 between two variants. In the
\emph{rating} variant a majority of reviews share one star rating and 3 outliers carry a different
one (\emph{``find the outliers with the least common rating''}); in the \emph{category} variant the
majority share one product category and 3 outliers come from another (\emph{``find the outliers
from the least common category''}), with categories taken from the corpus's own product-category
labels. Ground truth comes from the structured metadata --- the star-rating field and the category
label --- and never from LLM labelling; since both fields are stripped from the visible text, the
model has to read the review content to recover them.

Because the label vocabulary here is \emph{fixed and small}, one labelling pass over the corpus
suffices and the task is \ctc{N} --- which is why it groups with Oolong rather than with the
Wikipedia outlier variants, whose category set is open. That contrast is the point of including
it. It also weakens the topical-similarity prior the Wikipedia variant enjoys, since reviews
sharing a rating often read alike, which makes it the harder of the two to learn without a
natural-language chain of thought. Reviews are short ($\sim$80 tokens), making this the densest
task in the suite in documents per token; we use 5{,}000 training and 500 evaluation examples.
\end{claudeblock}

\textbf{Grouping (OpenAlex).} \emph{Pairwise-F1 / ARI.} Given $N$ scientific abstracts and an integer $k$, partition them into $k$ groups. Gold partitions come from the OpenAlex \citep{Priem2022OpenAlexAF} concept hierarchy: we sample only papers annotated down to level L3, then per example pick a level $L \in \{L0,\ldots,L3\}$ and sample $k$ distinct concept values at that level. The level sets the granularity --- L0 gives 2--3 coarse groups, L3 gives 6--10 fine ones --- and, crucially, \textbf{the query states $k$ but never states the axis}, so the model must infer from the corpus and $k$ what granularity is being asked for. This makes grouping harder than outlier detection: the number of categories is a per-example variable rather than something to be discovered as ``one minority''. The instruction also asks for a short topic label per group before its IDs; only the IDs are scored, but we find that naming the group improves the partition itself.


\textbf{Cross-corpus absence (xabsence).} \emph{Set-F1.} Two corpora $A$ and $B$ share one context. Almost every document in one has a twin in the other; exactly $k$ are unmatched, and the model names those orphans.
\begin{claudeblock}
The suite uses the \textbf{exact-copy} variant: a twin is a byte-identical copy of its partner, and
documents are full PubMed abstracts rather than single sentences.

This replaces an earlier paraphrase variant, and the reason is a leak worth recording. In that
version a twin was a low-overlap LLM paraphrase, so corpus $B$ consisted almost entirely of
LLM-written text --- while an orphan inserted into $B$ was in original claim form. The orphan could
therefore be spotted from the style of that \emph{one} document, with no cross-corpus comparison at
all, and the asymmetry shows up plainly in the scores: a trained 4B model reached 0.98 recall on
$B$-side orphans against 0.08 on $A$-side ones. Making the twin an exact copy removes every
per-document cue by construction --- a matched document and an orphan are drawn from the same
distribution and are individually indistinguishable --- so the orphan can only be found by
comparing documents across the two corpora, which is the \ctc{N^2} computation the task is meant to
require. Using full abstracts also makes each document self-contained, at the cost that a
token-budgeted rung holds far fewer documents than the sentence variant did.
\end{claudeblock}


\textbf{Query-document matching (qdmatch).} \emph{Pair set-F1.} The retrieval analogue of contradiction, and our sharpest test that the \emph{same} substrate changes class when the query structure changes. An example interleaves $M$ queries and $N$ documents under one shared numbering; only $k$ of the queries are relevant, and each of those maps to all of its gold documents. Pools are drawn so that relevant queries, distractor queries (gold withheld) and distractor documents (query withheld) are disjoint, which guarantees the only true matches are the planted ones. The model must find the sparse gold pairs among $M \times N$ combinations. We build three versions from three retrieval sources --- NQ (1 gold/query), HotpotQA bridge (2 golds/query), and \claude{FiQA (median 2 golds/query, sparsely judged)} --- so difficulty from the underlying retrieval problem varies while the search structure is fixed.
% ⚠ CORRECTED: this said OBLIQ. The suite's three qdmatch tasks are qdmatch_nq, qdmatch_hpqa and
% qdmatch_fiqa -- there is no qdmatch_obliq anywhere in the data or the result tables, and the
% CTC-QDMATCH figure in the main paper plots NQ / HotpotQA / FiQA. A layout ablation presents the items either as a query block followed by a document block, or fully interleaved.

\subsubsection{\texorpdfstring{\ctc{N^2}}{OT(N2)} tasks: all-pairs search}\autonote{8 tasks, 6 new. Contradiction and
reorder are edited old blurbs (contradiction now names both perturbation modes and the
one-LLM-sentence-per-pair property). Metric wording fixed: contradiction was labelled ``EM'' only,
but is reported as pair set-F1 with EM alongside.}

% The oracle operation here relates two arbitrary documents, $a(d_i, d_j)$, and any of the $\binom{N}{2}$ pairs could be gold, so a naive solution is a nested loop. This is the regime we expect \cmc{N} attention to fail in, and most of these tasks do not exist in prior benchmarks.

\textbf{Contradiction search (PubMed).} \emph{Pair set-F1; EM = did the model find all gold pairs.} Each document is a single claim sentence from a PubMed abstract \citep{Jin2019PubMedQAAD}. Among $N$ documents, $K{=}3$ gold pairs $(d, d')$ are hidden: $d$ is a real PubMed sentence and $d'$ is a minimally-perturbed version, produced by an LLM, that factually contradicts it --- either an obvious polarity flip (``\emph{reduced} MI risk by 23\%'' $\to$ ``\emph{increased} MI risk by 23\%'') or a subtle numeric/scope edit (``30\% reduction'' $\to$ ``3\% reduction''); we mix both. \textbf{Only one sentence per pair is LLM-written}, which avoids the usual pitfall where the planted item is identifiable by writing style alone. Candidate sentences are filtered to be self-contained and claim-bearing (no leading anaphora, no dangling citations), and fillers come from abstracts disjoint from every gold-source abstract in that example.


\textbf{String matching (strmatch).} \emph{Pair set-F1.} A synthetic-substrate control for the same computation: $N$ strings of $L$ words drawn from a Wikipedia vocabulary, and the task is to find every pair sharing a contiguous run of $\geq k$ words. Gold pairs share exactly one $k$-word run; hard negatives share a $(k{-}1)$-word run, so the model must count run length rather than notice overlap. Every non-planted word is globally unique in the example, which makes the gold set provably exhaustive. This task has no semantic content at all, so the only thing that scales is the number of comparisons.



\textbf{Reordering.} \emph{Kendall-$\tau$.} $N$ consecutive segments of a single Project Gutenberg book, presented in random order, with the task of recovering the original ordering as a permutation of document IDs --- a harder version of the reconstruction tasks in \citet{Chen2026LongBenchPA}. Chapter headings are stripped so the model cannot shortcut via ``Chapter 3''. Because sorting requires comparisons between arbitrary pairs of segments, we label it \ctc{N^2}. A Wikipedia variant (consecutive 100-word chunks from one article) exists as a harder, less narrative-sequential control.

\subsubsection{Beyond \texorpdfstring{\ctc{N^2}}{OT(N2)}}
% The \autonote here named cycle and groups-of-4 alongside textgroups. Both of those were cut from
% the roster and their descriptions are already commented out below, so the note pointed at text
% that no longer exists. Textgroups is the only surviving task in this class.

% Three tasks require joint consideration of triples or longer chains, which no pairwise sweep can produce. They are deliberately short in tokens: the point is that difficulty scales with the number of \emph{items}, not the raw context length.

% \textbf{Cycle detection.} \emph{Group set-F1.} $N$ comparative claims (``\emph{A scored higher than B}'') over named entities, all with the same predicate, so the corpus defines one directed graph. A strict order admits no directed cycle, so any cycle is an inconsistency; the task is to name the claims forming each one. The construction gives every entity a single global rank and makes every background edge point forward in that rank, so the background graph is a DAG and each planted cycle is provably the unique one closing its backward edge. Cycle entities participate in background edges like any other entity, so gold claims are not identifiable by frequency --- an earlier version failed this and was rebuilt. Finding a length-$L$ cycle naively is \ctc{N^L}.

% \textbf{Groups-of-four.} \emph{Group set-F1.} $N$ arithmetic expressions; find every group of $G{=}4$ whose values are all within $X$ of each other. Distractors are drawn to be in-distribution --- a distractor may lie within $X$ of up to $G{-}2$ others, so incidental close pairs and triples occur --- but the construction never lets a run of $G$ mutually-close values form outside the planted groups, so the gold group stays unique while close pairs stop being a giveaway.

\textbf{Textgroups.} \emph{Group set-F1.} The textual analogue: $N$ short passages, each carrying a feature the model must read off the prose (how many nouns / verbs / adjectives it contains, or how often a chosen connective appears), and the task is to find every triple whose feature values sum to a target $T$. The lexicons are closed and pairwise disjoint so the count is unambiguous, and sentence structure is varied so the feature is decoupled from passage length. Because $G{=}3$, the construction can both plant exactly $K$ triples and verify by brute force that no others exist. This is the one task in the suite where the per-document oracle operation is itself nontrivial (counting) \emph{and} the combination is super-quadratic.
